\section{Related Work}
\label{sec:related}

\paragraph{Empirical critiques of algorithmic audit regimes.} A growing body of FAccT scholarship documents that existing algorithmic audit regimes fail to deliver the accountability they promise, largely through empirical and doctrinal analysis. Wright et al.~\cite{wright2024null} study 391 NYC employers' Local Law 144 compliance and introduce \emph{null compliance}: non-compliance that cannot be ascertained because the statute grants employers discretion over what is audited and disclosed. Groves et al.~\cite{groves2024auditing} similarly document structural weaknesses in the NYC regime's implementation. Neither formalizes a dynamic model or predicts investment changing systematically around certification. Power's account of auditing as institutional ``comfort rather than treatment''~\cite{power1997audit} is a qualitative diagnosis our model gives specific, falsifiable, dynamic form.

\paragraph{Enforcement economics.} That compliance investment responds to monitoring probability organizes the economics of law enforcement, from Becker's deterrence analysis~\cite{becker1968crime} through the inspection-game literature~\cite{mookherjee1992monitoring}. We adapt this to point-in-time algorithmic certification, where the discrete uncertified-to-certified transition is itself the event that changes perceived monitoring probability, deriving the kink corollary as a consequence of this structure.

\paragraph{Post-certification compliance relaxation elsewhere.} That passing an audit lowers subsequent compliance effort is documented, but not formalized dynamically, elsewhere. Kasper and Rablen~\cite{kasper2023tax} isolate a tax-compliance ``bomb-crater effect'': an audit lowering perceived re-audit probability causes compliance to fall, the closest empirical analogue to our mechanism. Gray and Mendeloff~\cite{gray2004declining} document OSHA inspections' injury-reducing effect attenuating from 19\% (1979--85) to statistically insignificant (1992--98). Both are empirical rather than formal.

\paragraph{Dynamic and repeated inspection.} Varas, Marinovic, and Skrzypacz~\cite{varas2020random} solve a continuous-time model where reputation depreciates between inspections and is re-earned, the closest scaffold for lapse-and-re-earn certification (Section~\ref{sec:recurring}). Che, Kim, and Mierendorff~\cite{che2024datadriven} extend this to a regulator learning compliance from noisy outcomes over time; Section~\ref{sec:imperfectobs} relaxes our own perfect-observation assumption more simply, finding cycling eliminated only once noise exceeds a threshold — a qualitative connection, where it holds, to their finding that imperfect information can sustain deterrence a fully-resolved state cannot. Fernandes and Phelan~\cite{fernandes2000recursive} and Halac and Prat~\cite{halac2016managerial} instead solve the fully forward-looking dynamic program against a proper continuation value; we adopt the myopic simplification for tractability (Section~\ref{sec:discussion}). Ball and Knoepfle~\cite{ball2023timing} show random inspection \emph{timing} outperforms predictable timing against decay — a lever distinct from our \emph{detection}-noise result, since our model normalizes cadence to one period.

\paragraph{Regulator-side mechanism design.} Macho-Stadler and P\'{e}rez-Castrillo~\cite{machostadler2006optimal} solve optimal emissions-audit allocation across strategically under-reporting firms; Huang and Zhu~\cite{huang2024zetar} give a Stackelberg information-design framework for privately-known types. Both treat non-compliance as static, without a decaying or lapse-and-re-earn state. None combines a certification event that shifts beliefs, a provable non-monotone effort trajectory, and a regulator anticipating it — what Section~\ref{sec:regulator} takes a first step toward, with the joint design problem under heterogeneity left open.

\paragraph{Algorithmic audit-gaming, institutional decoupling, and multiplicity.} Garcia Bourr\'{e}e et al.~\cite{garciabourree2025robust} and Basu and Das~\cite{basu2025fairgame} ask whether a \emph{fixed} audit can be certified manipulation-proof; our result is structurally different, since even a manipulation-proof, honestly-measured predicate can produce collapse through legitimate investment choices, not measurement gaming. Our claim that certification can decouple from the property it certifies echoes organizational sociology's ``myth and ceremony''~\cite{meyer1977institutionalized,edelman1992legal,bromley2012smoke}, which diagnoses decoupling generally where we contribute a specific formal mechanism. Model-multiplicity work~\cite{marx2020predictive,black2022model} shows a \emph{cross-sectional} indeterminacy in which model is certified, distinct from our \emph{temporal} claim about one model's quality evolving post-certification; De Toni et al.~\cite{detoni2025time} show individual recourse validity degrades under drift, a parallel temporal-invalidation result at the level of individual guarantees rather than firm-level incentives~\cite{lyons2021conceptualising}.
